\subsection{Boundary types and the scope of the covering criterion}

\begin{statement}{E.27}{Proved}{The four boundary types and their actual comparison maps}
Over an algebraically closed field of characteristic zero the geometric
fundamental groups of types $(0,0),(0,1),(0,2),(1,0)$ are, respectively,
\[
 1,\qquad1,\qquad\widehat{\mathbb Z}(1),\qquad
 T E\simeq\widehat{\mathbb Z}^{,2},
\]
where a chosen origin gives the elliptic curve $E$ in the last case.
The last two groups are their own centers and their own abelian and
metabelian quotients. The assertions about nontrivial comparison kernels
in E.7--E.9 have the stated rank-two free-group source. For the discrete
group $F_1=\mathbb Z$ and its subgroup $d\mathbb Z$, the abelian,
metabelian and fixed nilpotent comparisons to their correctly descended
kernels are isomorphisms for every $d$. The rational completion map
is multiplication by $d$ on $\mathbb Q$, hence is also an isomorphism.
For $d=p^a m$ with $(m,p)=1$, the pro-$p$ comparison is the isomorphism
\[
 \mathbb Z_p\longrightarrow p^a\mathbb Z_p,\qquad u\longmapsto du.
\]
The same statements hold for a finite-index lattice in $\mathbb Z^2$,
with its inclusion matrix in place of $d$.

For the corresponding nontrivial finite covers of $B\mathbb Z$ or
$B\mathbb Z^2$, intrinsic-fibration membership for the five localizers
of E.5 still holds only for the $p$-group covers in the pro-$p$ case.
Thus the equivalence between the two criteria in E.10 is proved for its
specified source $F_2$; the displayed lattice comparisons compute what
happens at these boundary sources.
\end{statement}
\begin{proof}
For the projective line a connected degree-$d$ finite etale cover has
$2g-2=-2d$ by Riemann--Hurwitz, forcing $d=1$. For the affine line,
extend a finite etale cover to smooth projective curves. Ramification
can occur only over infinity, and its total contribution is at most
$d-1$. Hence $2g-2\le-d-1$, again forcing $d=1$. For the twice-punctured
line, ramification is allowed at two points. The same formula forces
$g=0$ and total ramification over each of them. Choosing coordinates
at their unique preimages makes the covering $z\mapsto z^d$.
The cofinal inverse system of these maps has deck groups $\mu_d$,
which proves the asserted Tate-module description and its action.

A finite etale cover of a genus-one curve has genus one by the same
formula. Choose a point over the fixed origin. The resulting pointed
map of elliptic curves is an isogeny: its map to the target Jacobian
is the pointed curve map under the degree-zero divisor identification.
Its kernel is finite; its dual isogeny shows that a multiplication
map dominates it. The multiplication coverings are therefore cofinal,
and their inverse deck system is $TE$. In characteristic zero
$E[n]\simeq(\mathbb Z/n)^2$, giving the displayed geometric group.
For the cubic over $\mathbb Q$ the specified origin is a rational
coordinate cusp; E.24 computes its required arithmetic abelian
cohomology. We impose no rational-point conclusion on a genus-one
torsor without such an origin.

For the rank-one discrete group, every verbal commutator subgroup is
zero, so the covering inclusion itself is exactly the displayed
kernel map. Rationalization gives multiplication by $d$.
In $\mathbb Z_p$, multiplication by $m$ is a unit and multiplication
by $p^a$ maps bijectively onto the closed subgroup $p^a\mathbb Z_p$.
For a rank-two lattice, integral Smith normal form gives the same
calculation in two coordinates. These arguments compute both the
image and zero kernel, rather than applying the free-group rank
formula outside its hypotheses. The program \texttt{e\_boundary.py}
checks the rank-one formulas for $d\le24$ and the inverse unit
coordinates modulo $p^k$, for $p=2,3,5,7$ and $k\le4$.

Intrinsic-fibration membership is preserved by pullback, directly
from its universal base-change definition. Pull a nontrivial lattice
cover back by $BF_2\to B\mathbb Z$ or $BF_2\to B\mathbb Z^2$,
using a surjection onto the lattice. Its monodromy is still the
given finite quotient, so E.10 excludes the pulled-back cover in
all the claimed cases. The general $p$-group sufficiency in E.10
proves the other cases. This supplies the comparison between the
two criteria at every boundary source under discussion.

Real descent also retains its arithmetic data. The plus conic
$X_2:x_0^2+x_1^2+x_\infty^2=0$ has no real point, but its
geometric projective fundamental group is trivial, so its arithmetic
fundamental group maps isomorphically to $G_{\mathbb Q}$ and has
its unique section. Deleting the coordinate divisors gives the
nontrivial real obstruction treated by E.25. The exact open-to-proper
quotient kills that obstruction. The elliptic and multiplicative
centers above, the retained central $C_2$ in A.14, and the real
orientation-reversing extension in E.18 are consequently part of
the calculations rather than omitted exceptional cases.
\end{proof}
